\iftoggle{paper}
{
This document presents the thesis for obtaining the Master's degree in Computer Science: Emphasis in Computer Science at the Tecnológico de Costa Rica.
The topic is centered in the idea of the use of machine-learning techniques to learn patterns from acoustic echoes produced inside small rooms and estimate the location of objects in three-dimensional space. Two neural-network architectures are employed for this task: a Convolutional Neural Network (CNN), which, although primarily designed for image processing, is commonly used in the literature for acoustic inference; and an Echo State Network (ESN), a memory- and computationally efficient architecture with the potential to process acoustic signals more directly. The effect of beamforming is also investigated as a possible solution to multipath acoustic interference. The network designs, experimental setup, and methodology for measuring, processing, and evaluating the results are presented to assess the potential of audible-frequency sound echoes for locating objects in an enclosed room.
}
{}
